%% LyX 2.4.4 created this file.  For more info, see https://www.lyx.org/.
%% Do not edit unless you really know what you are doing.
\documentclass[english]{article}
\usepackage[T1]{fontenc}
\usepackage[latin9]{inputenc}
\usepackage{xcolor}
\usepackage{booktabs}
\usepackage{calc}
\usepackage{units}
\usepackage{enumitem}
\usepackage{amsmath}
\usepackage{amsthm}
\usepackage{cancel}
\usepackage{geometry}
\geometry{verbose,tmargin=1in,bmargin=1in,lmargin=1in,rmargin=1in}

\makeatletter

%%%%%%%%%%%%%%%%%%%%%%%%%%%%%% LyX specific LaTeX commands.
%% Because html converters don't know tabularnewline
\providecommand{\tabularnewline}{\\}

%%%%%%%%%%%%%%%%%%%%%%%%%%%%%% Textclass specific LaTeX commands.
\numberwithin{equation}{section}
\numberwithin{figure}{section}
\newlength{\lyxlabelwidth}      % auxiliary length 

%%%%%%%%%%%%%%%%%%%%%%%%%%%%%% User specified LaTeX commands.
\usepackage{fancyhdr}
\usepackage{lastpage}
\pagestyle{fancy}
\usepackage{enumitem}
\usepackage{ifthen}
%sets math fonts to be more readable
\everymath=\expandafter{\the\everymath\displaystyle}
%\DeclareMathSizes{10}{10}{10}{10}
\usepackage{multicol}

\usepackage{tikz}
\usetikzlibrary{plotmarks}
\usepackage{pgfplots}
\pgfplotsset{compat=newest}

\usepackage{calc}
\usepackage[nomessages]{fp}

\usepackage{advdate}    % Advancing/saving dates
\usepackage{datetime}   % Dates formatting
\usepackage{datenumber} % Counters for dates
% Added by lyx2lyx
\setlength{\parskip}{\smallskipamount}
\setlength{\parindent}{0pt}

\makeatother

\usepackage{babel}
\begin{document}
\renewcommand{\author}{Dr.\,James Ashe}

\newcommand{\classNum}{132}

\newcommand{\classSec}{C and K}

\newcommand{\nameType}{Solutions}

\newcommand{\examType}{Worksheet 5.1 and 5.2}

\renewcommand{\date}{\formatdate{18}{3}{2014}}

%%First page's header%%%%%%%%%%%%%%%%%%%%%%
\fancypagestyle{firststyle} %%header settings for first pagr
{
	\fancyhead[L]{MTH \classNum{}(\classSec{}) \author{}\\
	\textbf{\examType{}} \date{}}
	\fancyhead[C]{\textbf{\nameType}}
	\setlength{\headsep}{14pt}
}
%%%%%%%%%%%%%%%%%%%%%%%%%%%%%%%%%%%%%%%%%%

%%Next page's header%%%%%%%%%%%%%%%%%%%%%%%
\setlist{nolistsep}
\lhead{\classNum{}(MTH \classSec{})} 
\chead{\textbf{\nameType}} 
\cfoot{\thepage{}~of~\pageref{LastPage}} 
\setlength{\headsep}{5pt} 
%%%%%%%%%%%%%%%%%%%%%%%%%%%%%%%%%%%%%%%%%%%

%%Various settings; see the preamble for other settings%%%%%
%%\setenumerate[0]{leftmargin=12pt,itemindent=0pt,labelwidth=0pt}
\renewcommand{\labelenumi}{\textbf{\arabic{enumi}.}}
\renewcommand{\labelenumii}{\textbf{\alph{enumii}.}}
\thispagestyle{firststyle}
%%%%%%%%%%%%%%%%%%%%%%%%%%%%%%%%%%%%%%%%%%
\newcommand{\xmax}{10}
\newcommand{\xmin}{-10}
\newcommand{\xstep}{0} %must be xmin+n
\newcommand{\ymax}{10}
\newcommand{\ymin}{-10}
\newcommand{\ystep}{0} %must be ymin+n

\newcommand{\xspace}{\vspace{.5cm}}

\textbf{Formulas 5.1}

\noindent\begin{minipage}[t]{1\columnwidth}%
\begin{tabular}{llll}
$P=$principal/present value & $r=$annual interest rate & $m=$compoundings per year & $n=tm$\tabularnewline
$A=$future/maturity value & $t=$number of years & $i=\nicefrac{r}{m}=$rate per period & $r_{E}=$effective rate\tabularnewline
 &  &  & \tabularnewline
\end{tabular}%
\end{minipage}

\noindent\fcolorbox{black}{white}{\begin{minipage}[t]{1\columnwidth - 2\fboxsep - 2\fboxrule}%
\begin{tabular}{lclcll}
\textbf{Simple Interest} &  & \textbf{Compound Interest} &  & \textbf{Continuous Compounding} & \tabularnewline
\cmidrule{1-5}
$A=P(1+rt)$ &  & $A=P(1+i)^{n}$ &  & $A=Pe^{rt}$ & Future\tabularnewline
$P=\frac{A}{1+rt}$ &  & $P=A(1+i)^{-n}$ &  & $P=Ae^{-rt}$ & Present\tabularnewline
 &  & $r_{E}=\left(1+\frac{r}{m}\right)^{m}-1$ &  & $r_{E}=e^{r}-1$ & Effective\tabularnewline
\end{tabular}%
\end{minipage}}

\textbf{\xspace}

\textbf{Formulas 5.2}

\noindent\begin{minipage}[t]{1\columnwidth}%
\begin{tabular}{lll}
$S=$future value &  & $n=$number of periods\tabularnewline
$R=$periodic payment &  & $i=$interest rate per period\tabularnewline
 &  & \tabularnewline
\end{tabular}%
\end{minipage}

\noindent\fcolorbox{black}{white}{\begin{minipage}[t]{1\columnwidth - 2\fboxsep - 2\fboxrule}%
\begin{tabular}{lcl}
\textbf{Future Value of an Ordinary Annuity} &  & \textbf{Sinking Fund Payment}\tabularnewline
\midrule
$S=R\left[\frac{(1+i)^{n}-1}{i}\right]$ &  & $R=\frac{S\cdot i}{(1+i)^{n}-1}$\tabularnewline
\end{tabular}%
\end{minipage}}

\xspace

\textbf{Find the simple interest AND the future value of the loan. }
\begin{enumerate}[resume]
\item $\$1500$ at 10\% for 10 months.

\begin{enumerate}[resume]
\item [FV:]$A=1500(1+.1\frac{10}{12})=\$1625$
\item [I:]$1625-1500=\$125$\vfill
\end{enumerate}
\item \$3000 at 8\% for 9 months.\vfill
\item \$47,400 at 15.75\% for 157 days. Assume 360 days in a year and 30
days in a month. \vfill
\end{enumerate}
\textbf{Find the compound amount for the given deposit and the amount
of interest earned.}
\begin{enumerate}[resume]
\item \$1100 at 7.3\% compounded semiannually for 10 years.

\begin{enumerate}[resume]
\item []$A=1100(1+\frac{.073}{2})^{10\cdot2}\approx\$2253.09$\vfill
\end{enumerate}
\item \$1100 at 7.3\% compounded monthly for 52 months.\vfill
\item \$1100 at 7.3\% compounded continuously for 4 years.\vfill
\end{enumerate}
\textbf{\newpage}

\textbf{Find the amount that should be invested now to accumulate
the following amount, if the money is compounded as indicated.}
\begin{enumerate}[resume]
\item \$5500 at 15.75\% compounded monthly for 39 months. 

\begin{enumerate}[resume]
\item []$t=\frac{39}{12}=3.25$\\
$P=5500(1+\frac{.1575}{12})^{-3.25\cdot12}\approx\$3307.53$
\end{enumerate}
\item \$6000 at $3.5\%$ compounded quarterly for 3 years.\vfill
\item \$6000 at $3.5\%$ compounded continuously for 3 years. \vfill
\end{enumerate}
\textbf{Solve}
\begin{enumerate}[resume]
\item A company will need \$100,000 in 9 years for a new building. To meet
this goal the CFO will invest money into an account paying 9\% annual
interest rate compounded quarterly. How much money needs to be invested
into the account?

\begin{enumerate}[resume]
\item []$100,000=(1+\frac{.09}{4})^{4\cdot9}$
\end{enumerate}
\item You win the lottery. You are offered \$13,00 now or \$20,200 in 5
years. If the money can invested at 5.5\% compounded annually, which
prize will be worth more in 5 years?\vfill
\end{enumerate}
\textbf{\newpage}

\textbf{Find the future value of each ordinary annuity, if payments
are made and interest is compounded as given. Then determine how much
of this value is from interest.}
\begin{enumerate}[resume]
\item $R=\$9,200$ (payments); $10\%$ interest compounded semiannually
for 7 years.\vfill
\item Payments of $\$800$ into an account paying 6.51\% interest compounded
quarterly for 12 years. \vfill
\end{enumerate}
\textbf{Determine the interest rate needed to accumulate the following
amounts in a sinking fund, with payments as given. }
\begin{enumerate}[resume]
\item Accumulate \$56,000, monthly payments of \$300 over 12 years.\vfill
\item Accumulate \$120,000, monthly payments of \$500 over 15 years. \vfill
\end{enumerate}
\textbf{Find the periodic payment that will amount to each given sum
under the given conditions. }
\begin{enumerate}[resume]
\item $S=\$100,000$; interest is 5\% compounded annually; payments are
made at the end of each year for 12 years. \vfill
\item $S=\$200,000$; interest is 5\% compounded annually; payments are
made at the end of each year for 12 years. \vfill
\end{enumerate}

\end{document}
